\documentclass[11pt,a4paper]{article}
\usepackage[T1]{fontenc}
\usepackage[utf8]{inputenc}
\usepackage{lmodern,microtype}
\usepackage[margin=25mm,headheight=14pt]{geometry}
\usepackage{amsmath,amssymb,amsthm,mathtools,mathrsfs}
\usepackage{booktabs,array,longtable,enumitem}
\usepackage[dvipsnames]{xcolor}
\usepackage{fancyhdr,xurl}
\usepackage[colorlinks=true,linkcolor=MidnightBlue,citecolor=MidnightBlue,urlcolor=MidnightBlue,unicode]{hyperref}
\usepackage{bookmark}
\hypersetup{pdftitle={Natural Boundaries Survive Nonlinear Feedback},pdfauthor={Research article prepared with ChatGPT for Vladimir Reshetnikov},pdfsubject={Partial theta level sets, transseries, fine versus angular Borel summation}}
\pagestyle{fancy}\fancyhf{}
\fancyhead[L]{\small Natural boundaries and nonlinear feedback}
\fancyhead[R]{\small Research article}
\fancyfoot[C]{\thepage}
\renewcommand{\headrulewidth}{0.3pt}
\setlength{\parskip}{3pt}
\setlength{\emergencystretch}{3em}
\setlist{leftmargin=1.8em,itemsep=4pt,topsep=5pt}
\numberwithin{equation}{section}
\newtheorem{theorem}{Theorem}[section]
\newtheorem{proposition}[theorem]{Proposition}
\newtheorem{lemma}[theorem]{Lemma}
\newtheorem{corollary}[theorem]{Corollary}
\theoremstyle{definition}\newtheorem{definition}[theorem]{Definition}
\newtheorem{question}[theorem]{Question}
\theoremstyle{remark}\newtheorem{remark}[theorem]{Remark}
% ed. (2026-09-29): unnumbered environment for ProveIt editorial notes; it does not shift any numbering.
\theoremstyle{definition}
\newtheorem*{ednoteinner}{Editorial note (ProveIt, 2026-09-29)}
\newenvironment{ednote}{\begin{ednoteinner}\sloppy\Urlmuskip=0mu plus 1mu\relax}{\end{ednoteinner}}
\newcommand{\C}{\mathbb C}\newcommand{\R}{\mathbb R}\newcommand{\N}{\mathbb N}\newcommand{\Z}{\mathbb Z}
\newcommand{\Uhat}{\widehat U}\newcommand{\Qhat}{\widehat Q}
\newcommand{\Rea}{\operatorname{Re}}\newcommand{\Ima}{\operatorname{Im}}
\newcommand{\dd}{\,\mathrm d}\newcommand{\e}{\mathrm e}
\newcommand{\B}{\mathcal B_0}\newcommand{\Hh}{\mathcal H}
\newcommand{\abs}[1]{\lvert #1\rvert}
\newcommand{\norm}[1]{\lVert #1\rVert}
\newcommand{\co}[2]{[#1^{#2}]}
\newcommand{\inv}{\langle-1\rangle}
\begin{document}
% ed. (2026-09-29): no page anchors on the title page, which the titlepage environment numbers 1 like the following page (removes a duplicate-destination warning).
\hypersetup{pageanchor=false}
\begin{titlepage}
\centering\vspace*{10mm}
{\small\scshape Research in transseries and analytic continuation\par}
\vspace{9mm}
{\LARGE\bfseries Natural Boundaries Survive\par Nonlinear Feedback\par}
\vspace{7mm}
{\Large Moving-parameter partial theta series,\par sharp half-plane realization,\par and the failure of angular Borel summation\par}
\vspace{11mm}
{\large Prepared for Vladimir Reshetnikov\par}
\vspace{3mm}29 September 2026
\vspace{10mm}
\begin{minipage}{0.94\textwidth}
\begin{abstract}
For an integer $d\ge2$, consider the formally well-defined equation
\[
 \Uhat_d(q)=\sum_{j\ge1}q^j\exp\!\bigl(j^d\Uhat_d(q)\bigr)
\]
and its compositional inverse $\Qhat_d$. We prove that the canonical actual
inverse $Q_d$, holomorphic in a sufficiently small left half-disc and smooth
up to its diameter, has that entire imaginary diameter as a natural boundary.
The obstruction survives analytic changes of the moving parameter, analytic
changes of the level set, and finite modifications of the feedback kernel.
The main tool is an explicit large-order theorem: periodic geometric moments
retain factorial growth after substitution along an arbitrary holomorphic
curve. A mean-square argument treats tied nearest poles without assuming a
unique dominant singularity.

% ed. (2026-09-29): at d = 2 this answer duplicates an earlier-filed package; see the editorial notes in Section 1.1 and after Theorem thm:quadratic-main.
For $d=2$, this resolves the angular-continuation question left open in the
repository's negative-ray summation article: both inverse and forward formal
series are finely Borel summable on the negative ray, but neither is
Borel summable in any open angular neighborhood of that direction. In
particular, an entire Borel transform does not supply an angular summation
sector. We also give convergent exponential blocks, uniform closed-half-disc
remainders for every $d\ge2$, an explicit optimized error scale, finite
verification programs, and further research questions.
\end{abstract}
\end{minipage}
\vfill
\begin{minipage}{0.94\textwidth}\small
\textbf{Proof status.} The new assertions are accompanied by conventional
mathematical proofs. They have not been checked in Lean or independently peer
reviewed. The contribution is a model-specific natural-boundary and
non-summability theorem, not a claimed solution of a named longstanding
conjecture. The literature and repository comparison is targeted; global
priority is not asserted.
\end{minipage}
\end{titlepage}
\hypersetup{pageanchor=true}
\setcounter{tocdepth}{2}\tableofcontents\clearpage

\section{The question resolved and the contribution boundary}
\label{sec:context}
\subsection{The precise repository problem}
The ProveIt transseries documents distinguish formal coefficient calculations,
actual analytic realizations, remainder estimates, and summability. Those
levels must not be conflated. In the feedback family
\begin{equation}
 \Uhat(q)=\sum_{j\ge1}q^j\e^{\lambda_j\Uhat(q)},\qquad \lambda_j\ge0,
 \label{eq:general-feedback}
\end{equation}
several recent repository articles establish regularity classifications and
negative-direction realizations \cite{repo-reg,repo-sector,repo-fine}.
The most direct predecessor, \emph{Negative-Ray Summation of Countable
Exponential-Feedback Transseries}, proves fine negative-ray summability at the
threshold $\lambda_j=O((j\log(j+1))^2)$. It explicitly does not claim angular
summability or a maximal continuation domain \cite{repo-fine}.

This paper addresses the missing angular question for the quadratic model
$\lambda_j=j^2$, and proves a natural-boundary theorem for the larger family
$\lambda_j=j^d$, $d\in\N$, $d\ge2$. The answer in the quadratic case is
\emph{negative}: fine summation cannot be enlarged to an open angular interval
containing the negative direction. The reason is a natural boundary of the
actual inverse, not failure of the Borel germ to continue to finite points.
% ed. (2026-09-29): the d = 2 answer is not new to the repository.
\begin{ednote}
For $d=2$ this question had already been answered negatively, independently,
by an earlier-filed package that this article could not see \cite{ed:nbq}; see
the editorial note after Theorem~\ref{thm:quadratic-main}. The cases
$d\ge3$ are new to the repository.
\end{ednote}

The source snapshot used for the final comparison is
\begin{center}\small\texttt{ebd8344bca77d8352cd745b7df374618a290a029}.\end{center}
The principal files are under
\path{Analysis/Transseries/docs/series-and-transseries/}.
The source inventory and exact blob identifiers are recorded in the package's
provenance file. No repository files were changed. The canonical large volume
was not audited line by line, and no assertion about the entire repository's
formal proof status is needed here.

\subsection{Main statements}
Write
\begin{equation}
 \Theta_d(q,u)=\sum_{j\ge1}q^j\e^{j^d u},\qquad
 H_r^- =\{u:|u|<r,\ \Rea u<0\}.
 \label{eq:theta}
\end{equation}
For $|q|<1$, the series and all its mixed derivatives extend continuously to
$\Rea u=0$. The defining identity for the actual inverse is
\begin{equation}
 \Theta_d(Q_d(u),u)=u,\qquad Q_d(u)=u+O(u^2).
 \label{eq:inverse}
\end{equation}

\begin{theorem}[Natural boundary of the inverse and forward realization]
\label{thm:main}
For each integer $d\ge2$, there is $r>0$ such that:
\begin{enumerate}[label=(\roman*)]
\item Equation \eqref{eq:inverse} has a unique small solution $Q_d$ on
$H_r^-$. It is holomorphic there and $C^\infty$ up to the imaginary diameter,
including the origin.
\item No point of $\{it:|t|<r\}$ admits a holomorphic continuation of $Q_d$
through the diameter. Thus this whole boundary arc is a natural boundary.
\item After reducing $r$, $Q_d$ is one-to-one, and its inverse $U_d$ is
holomorphic on $\Omega_d=Q_d(H_r^-)$. The smooth arc
$\Gamma_d(t)=Q_d(it)$ is a natural boundary of $U_d$.
\item Uniformly as real $t\to0$,
\begin{equation}
 \Gamma_d(t)=it+2t^2-i\bigl(\tfrac92-2^d\bigr)t^3+O(t^4).
 \label{eq:boundary-shape-main}
\end{equation}
Neither canonical realization can be continued to a sector about the negative
axis with opening strictly greater than $\pi$.
\end{enumerate}
\end{theorem}
The phrase ``natural boundary'' here refers to the displayed boundary arc.
It is not a claim that the circular edge $|u|=r$ is singular or that the small
half-disc is a globally maximal domain.

\begin{theorem}[Fine but not angular in the quadratic case]
\label{thm:quadratic-main}
For $d=2$, the formal series $\Qhat_2$ and $\Uhat_2$ have entire integrated
Borel transforms and are finely Borel summable along the negative ray. Their
sums are $Q_2$ and $U_2$ near that ray. Neither formal series is $1$-summable in
an open angular interval containing direction $\pi$.

The same fine-versus-angular conclusion holds when finitely many of the
nonnegative slopes $j^2$ in \eqref{eq:general-feedback} are changed to other
finite nonnegative values.
\end{theorem}
Precise summability conventions are given in Section~\ref{sec:summability}.
Fine summability means a fixed-width Borel tube; angular summability means a
Borel sector with locally uniform exponential bounds. These are different
requirements.

% ed. (2026-09-29): the case d = 2 duplicates an earlier-filed package (batch 49), not visible from this article's snapshot.
\begin{ednote}
The case $d=2$ of Theorems~\ref{thm:main} and~\ref{thm:quadratic-main} was
proved independently in
\path{Natural_Boundaries_Quadratic_Exponential_Feedback/} \cite{ed:nbq} (its
Theorem~\texttt{thm:main}), filed earlier and not visible from this
article's snapshot, with the explicit radius $1/32$, the bound
$|Q'-1|<1/3$ and the further conclusion that the image arc is nowhere
real-analytic. Its Theorem~\texttt{thm:rigidity} is the case $d=2$ of
Theorem~\ref{thm:moving}, and its Theorem~\texttt{thm:rational-main} extends
the $d=2$ conclusions to rational amplitude series. The exact coefficients
at $d=2$ agree with that package's recorded ones (through degree 16). The cases $d\ge3$ answer
that package's question ``Higher-degree and nonpolynomial slopes'' for
integer $d$; noninteger powers remain open.
\end{ednote}

\subsection{What is new, inherited, and deliberately not claimed}
The principal new mechanism is the moving-curve moment theorem in
Section~\ref{sec:moving}. A fixed-parameter lacunary-series theorem alone does
not settle \eqref{eq:inverse}: its parameter $q$ changes with $u$, so one must
exclude cancellation by that motion. The resulting natural boundary and
angular obstruction are the proposed research contribution.
% ed. (2026-09-29): at d = 2 the moving-curve theorem and the obstruction are also in an earlier-filed package.
\begin{ednote}
At $d=2$ the moving-curve theorem, the natural boundary and the angular
obstruction are also in \cite{ed:nbq} (see the editorial note after
Theorem~\ref{thm:quadratic-main}); they are new to the repository for
$d\ge3$.
\end{ednote}

The coefficient/Gevrey classification and the positive fine-summability
result are earlier repository results. Sections~\ref{sec:blocks}--\ref{sec:summability}
recover the quadratic positive result in a specialized form to make the
contrast transparent. General Laplace summation and its nonlinear closure are
classical \cite{sokal,sauzin}. Partial-theta boundary asymptotics, rational
boundary points, and discrete Fourier transforms also have a substantial
literature \cite{hlss,folsom}. The new argument here concerns a geometric
parameter constrained by a nonlinear level-set equation, not a new discovery
of theta-series boundary phenomena.

We do not infer failure of resurgence merely from a natural boundary of the
summed function. Nor do we identify a Stokes constant from an upper remainder
bound. For $d>2$, the natural-boundary assertion is about the canonical
literal realization; no unsupported identification with every possible
higher-order summation procedure is made.

\section{Formal series and actual half-plane solutions}
\label{sec:realization}
\subsection{Formal existence is not analytic existence}
Let $\C[[q]]$ have its coefficientwise, or $q$-adic, topology. If $V,W$ have
zero constant term and agree modulo $q^n$, then
\[
 \sum_{j\ge1}q^j\e^{j^dV}
 \equiv\sum_{j\ge1}q^j\e^{j^dW}\pmod{q^{n+1}}.
\]
Each coefficient uses only finitely many $j$. Iteration beginning with zero
therefore yields a unique formal solution $\Uhat_d=q+O(q^2)$.
It has a unique compositional inverse $\Qhat_d=u+O(u^2)$, and substitution
in the formal identity gives
\[
 \sum_{j\ge1}\Qhat_d(u)^j\e^{j^d u}=u.
\]
No analytic implicit-function theorem at $(q,u)=(0,0)$ is being applied: the
kernel need not be holomorphic on a full product neighborhood of that point.
This distinction also matters when applying general Gevrey implicit-function
theorems such as \cite{nikolaev}.

Direct coefficient comparison gives
\begin{align}
 \Uhat_d(q)&=q+2q^2+\bigl(2^d+\tfrac72\bigr)q^3+O(q^4),\label{eq:first-U}\\
 \Qhat_d(u)&=u-2u^2+\bigl(\tfrac92-2^d\bigr)u^3+O(u^4).
 \label{eq:first-Q}
\end{align}
For $d=2$ the next coefficients are
\begin{align}
 \Uhat_2(q)&=q+2q^2+\tfrac{15}{2}q^3+\tfrac{107}{3}q^4
                 +\tfrac{1555}{8}q^5+O(q^6),\label{eq:quadratic-U}\\
 \Qhat_2(u)&=u-2u^2+\tfrac12u^3-\tfrac23u^4
                 -\tfrac{29}{8}u^5+O(u^6).
 \label{eq:quadratic-Q}
\end{align}
These finite identities are checked independently in the accompanying program.

\subsection{A uniform small root, including boundary values}
\begin{lemma}[Smooth half-plane implicit root]
\label{lem:root}
Let $f$ be holomorphic near zero, with $f(0)=0$. For every fixed $d\ge2$,
a sufficiently small closed left half-disc admits a unique small root
$q=Q_f(u)$ of
\begin{equation}
 \Theta_d(q,u)=f(u).
 \label{eq:forced}
\end{equation}
It is holomorphic in the open half-disc, $C^\infty$ up to its diameter, and
$Q_f(u)=f(u)+O(|u f(u)|+|f(u)|^2)$. The root vanishes exactly where $f$ does.
For $f(u)=u$, its boundary Taylor coefficients are those of $\Qhat_d$.
\end{lemma}
\begin{proof}
Separate the first term and use the fixed point
\[
 q=T_u(q):=\e^{-u}\left(f(u)-\sum_{j\ge2}q^j\e^{j^d u}\right).
\]
On $|q|\le R=1/8$, $\Rea u\le0$, the tail and its $q$-derivative are bounded by
\[
 \frac{R^2}{1-R},\qquad \frac1{(1-R)^2}-1.
\]
Choose $|u|\le1/16$ and then shrink the radius so that $|f(u)|\le R/4$.
The displayed bounds show that $T_u$ maps the closed $R$-disc strictly into
itself and has Lipschitz constant less than $1/2$. Uniform iteration gives the
root and holomorphic dependence for $\Rea u<0$. Comparison with $T_u(0)$
gives $|Q_f(u)|\le2\e^{|u|}|f(u)|$; reinsertion gives the stated expansion.

For each fixed pair of derivative orders, the mixed derivatives of the kernel
on a smaller $q$-disc are dominated by a convergent series of the form
$C\sum_{j\ge1}j^M R_1^j$, $R_1<1$. Thus every mixed derivative is continuous
up to $\Rea u=0$. The denominator $1-\partial_qT_u$ stays bounded away from
zero. Repeated implicit differentiation therefore extends all derivatives of
the root continuously to the diameter. This also proves the stated
$C^\infty$ regularity without assuming holomorphy across the diameter.

Since $\Theta_d(0,u)=0$, uniqueness implies the assertion about zeros.
Differentiation at zero, or finite Taylor comparison, identifies each finite
jet with the unique formal solution of the inverse equation.
\end{proof}
The $C^\infty$ conclusion does not assert convergence of the infinite Taylor
series. Smoothness at a boundary and holomorphic continuation through that
boundary are very different properties.

\section{Periodic moments and their pole geometry}
\label{sec:moments}
\subsection{A finite Fourier transform produces all relevant poles}
Let $(c_j)$ be a nonzero complex sequence of period $b$. Write
\begin{equation}
 \gamma_\ell=\frac1b\sum_{s=0}^{b-1}c_s\e^{-2\pi i\ell s/b},
 \qquad c_j=\sum_{\ell=0}^{b-1}\gamma_\ell\e^{2\pi i\ell j/b}.
 \label{eq:DFT}
\end{equation}
For $0<|q_0|<1$, define the geometric moments
\begin{equation}
 M_m=\sum_{j\ge1}c_jq_0^j j^m\qquad(m\ge0).
 \label{eq:moments}
\end{equation}
Choose one value of $\log q_0$; changing it only relabels the following poles.

\begin{lemma}[Moment singularities]
\label{lem:poles}
The exponential generating function of \eqref{eq:moments} is meromorphic:
\begin{equation}
 \mathcal M(z)=\sum_{m\ge0}\frac{M_m}{m!}z^m
 =\sum_{\ell=0}^{b-1}\gamma_\ell
  \frac{q_0\e^{z+2\pi i\ell/b}}{1-q_0\e^{z+2\pi i\ell/b}}.
 \label{eq:moment-function}
\end{equation}
Its nonremovable poles are precisely
\begin{equation}
 \mathcal A=\{-\log q_0-2\pi i\ell/b+2\pi i k:
              \gamma_\ell\ne0,\ k\in\Z\}.
 \label{eq:poles}
\end{equation}
At $a\in\mathcal A$, the principal part is $\gamma_a/(a-z)$, where
$\gamma_a$ is the corresponding Fourier coefficient. Let
\[
 \rho=\min_{a\in\mathcal A}|a|,\qquad
 J=\{a\in\mathcal A:|a|=\rho\}.
\]
Then $J$ has one or two elements. For some $R>\rho$,
\begin{equation}
 M_m=m!\sum_{a\in J}\frac{\gamma_a}{a^{m+1}}+O(m!R^{-m}).
 \label{eq:moment-asymp}
\end{equation}
In particular $|M_m|\le K m!\rho^{-m}$ for all $m$, with a fixed $K$.
\end{lemma}
\begin{proof}
For $z$ close to zero, absolute convergence permits summation of each
geometric series after \eqref{eq:DFT}; this gives
\eqref{eq:moment-function}. Distinct nonzero Fourier modes have disjoint pole
progressions, so their poles cannot cancel. All poles have real part
$\alpha=\log(1/|q_0|)>0$. Minimizing their modulus therefore minimizes the
absolute value of their imaginary part; at most one point of each sign can
attain that minimum. Moreover a nearest pole has imaginary part of modulus at
most $\pi$, because each individual progression has step $2\pi i$.

Remove the principal parts of the poles in $J$ and choose $R$ strictly between
$\rho$ and the next pole modulus. Cauchy's coefficient estimate on a circle of
radius $R$ gives \eqref{eq:moment-asymp}, with an arbitrarily slight reduction
of $R$ if needed. The final uniform bound follows, including finitely many
initial indices by enlarging $K$.
\end{proof}

\subsection{Why rational boundary points are enough}
At a rational imaginary point
\begin{equation}
 u_0=2\pi i a/b,\qquad a\in\Z,\quad b\in\N_{>0},
 \label{eq:rational-point}
\end{equation}
the sequence $c_j=\e^{j^d u_0}$ is periodic in $j$ and never zero. Indeed,
$(j+b)^d-j^d$ is divisible by $b$. No estimate for a Gauss sum is required
for the existence of a pole: a nonzero periodic sequence has at least one
nonzero Fourier coefficient.

For the quadratic case with $b$ odd and $\gcd(a,b)=1$, one additionally has
\begin{equation}
 |\gamma_\ell|=b^{-1/2}\qquad(0\le\ell<b).
 \label{eq:gauss}
\end{equation}
To check this, square the modulus of the unnormalized Fourier sum, put
$s=t+h$, and sum first over $t$. The inner sum is zero unless
$2ah=0\pmod b$, which forces $h=0$. Its remaining value is $b$.
Thus all $b$ pole progressions occur in this case. We use this identity only
for explicit diagnostics, not to bypass possible vanishing Fourier modes in
the general theorem.

\section{The moving-curve large-order theorem}
\label{sec:moving}
The central problem is whether an analytic choice of $q$ depending on $u$
can cancel all the boundary singularities. The next result rules this out.
It is a statement about prescribed holomorphic curves, not an unproved
asymptotic formula for the derivatives of the implicit solution itself.

Let $h$ be holomorphic near zero with $h(0)=0$, and let $h_1=h'(0)$. Consider
\begin{equation}
 F(v)=\sum_{j\ge1}c_jq_0^j\exp\!\bigl(j^d v+jh(v)\bigr),
 \qquad \Rea v\le0,
 \label{eq:moving-function}
\end{equation}
for sufficiently small $v$. This is the substitution
$q(v)=q_0\e^{h(v)}$. All derivatives extend continuously to zero from the
left half-plane. Write $f_n=F^{(n)}(0)/n!$ for these boundary coefficients.

\begin{theorem}[Persistence of the nearest-pole contribution]
\label{thm:moving}
With the notation of Lemma~\ref{lem:poles}, for $d=2$,
\begin{equation}
 f_n=\frac{(2n)!}{n!}\left(
       \sum_{a\in J}\frac{\gamma_a\e^{a h_1/2}}{a^{2n+1}}
       +o(\rho^{-2n})\right).
 \label{eq:moving-quadratic}
\end{equation}
For each integer $d\ge3$,
\begin{equation}
 f_n=\frac{(dn)!}{n!}\left(
       \sum_{a\in J}\frac{\gamma_a}{a^{dn+1}}
       +o(\rho^{-dn})\right).
 \label{eq:moving-higher}
\end{equation}
Thus a holomorphic motion changes the leading quadratic amplitude by a
nonzero factor. For $d\ge3$, it does not even change the leading amplitude.
\end{theorem}

\subsection{An exact coefficient identity}
Expand only a fixed derivative order. Absolute convergence of each moment
justifies the finite coefficient calculation
\begin{equation}
 f_n=\sum_{m=0}^{n}\sum_{k=0}^{m}
 \frac{[v^m]h(v)^k}{k!(n-m)!}\,M_{d(n-m)+k}.
 \label{eq:exact-moving}
\end{equation}
For $k=0$, only $m=0$ contributes. The term $(m,k)=(0,0)$ is $M_{dn}/n!$.
The identity does not assert convergence of a Taylor expansion of $F$.

Choose $H,R_h>0$ such that $|[v^\ell]h(v)|\le H R_h^{-\ell}$ for all
$\ell\ge1$. Then for $m\ge k\ge1$,
\begin{equation}
 |[v^m]h(v)^k|
 \le H^kR_h^{-m}\binom{m-1}{k-1}.
 \label{eq:h-bound}
\end{equation}
This follows by comparing $h$ with $H(v/R_h)/(1-v/R_h)$ coefficientwise in
absolute value.

\subsection{The factorial estimate that controls all indices}
For integers $0\le s\le N$,
\begin{equation}
 \frac{N!}{(N-s)!}\ge(N/\e)^s.
 \label{eq:falling}
\end{equation}
The geometric mean of the largest $s$ integers among $1,\dots,N$ is at least
the geometric mean of all $N$, and $N!\ge(N/\e)^N$.
Consequently
\begin{equation}
 \frac{n!}{(n-m)!}\frac{(dn-dm+k)!}{(dn)!}
 \le n^m\left(\frac{\e}{dn}\right)^{dm-k}.
 \label{eq:factorial-ratio}
\end{equation}
This bound is valid even when $m$ is close to $n$. An estimate proved only
for fixed $m$ would not justify the theorem.

Normalize the absolute value of a summand in \eqref{eq:exact-moving} by
\[
 L_n=\frac{(dn)!}{n!}\rho^{-dn}.
\]
Using \eqref{eq:h-bound}, \eqref{eq:factorial-ratio}, and the uniform moment
bound, and putting $m=k+r$, gives a bound of the form
\begin{equation}
 K\frac{X_n^k}{k!}\binom{k+r-1}{r}Y_n^r,
 \quad X_n=C_1n^{-(d-2)},\quad Y_n=C_2n^{-(d-1)}.
 \label{eq:global-bound}
\end{equation}
Here $C_1,C_2$ are fixed positive constants. For example one may take
$C_2=(\e\rho/d)^d/R_h$ and $C_1=C_2Hd/(\e\rho)$.
For $Y_n<1$,
\begin{equation}
 \sum_{k\ge1}\sum_{r\ge0}
 \frac{X_n^k}{k!}\binom{k+r-1}{r}Y_n^r
 =\exp\!\left(\frac{X_n}{1-Y_n}\right)-1.
 \label{eq:double-majorant}
\end{equation}
All extensions of finite sums to infinite sums here are made only for
nonnegative upper bounds.

\subsection{Completion of the asymptotic proof}
\begin{proof}[Proof of Theorem~\ref{thm:moving}]
When $d\ge3$, $X_n\to0$ and $Y_n\to0$. Equations
\eqref{eq:global-bound}--\eqref{eq:double-majorant} show that every term other
than $(0,0)$ contributes $o(L_n)$ in total. Apply
\eqref{eq:moment-asymp} to $M_{dn}/n!$ to obtain
\eqref{eq:moving-higher}.

For $d=2$, $X_n$ is constant and $Y_n=O(n^{-1})$. The part with $r\ge1$ is
bounded, after normalization, by
\[
 K\left\{\exp\!\left(\frac{X_n}{1-Y_n}\right)-\exp(X_n)\right\}
 =O(n^{-1}).
\]
It remains to evaluate the diagonal $m=k$, on which
$[v^k]h(v)^k=h_1^k$. A principal part $\gamma_a/(a-z)$ contributes
\[
 \frac{(2n)!}{n!}\frac{\gamma_a}{a^{2n+1}}
 \sum_{k=0}^n\frac{(a h_1)^k}{k!}
     \frac{n!}{(n-k)!}\frac{(2n-k)!}{(2n)!}.
\]
For each fixed $k$, the product of factorial ratios tends to $2^{-k}$.
The global bound supplies a summable majorant of the form $C^k/k!$.
Dominated convergence therefore makes the sum tend to $\e^{a h_1/2}$.

For the analytic remainder in \eqref{eq:moment-asymp}, each fixed diagonal
index gives $o(L_n)$, because $R>\rho$. The same summable majorant permits
dominated convergence over the diagonal. The off-diagonal part was already
$O(L_n/n)$. Summing over the finite set $J$ proves
\eqref{eq:moving-quadratic}.
\end{proof}

\begin{remark}[Scope of the amplitude formula]
The factor $\e^{a h_1/2}$ is nonzero regardless of the complex value of $h_1$.
This does not mean that two nearest-pole contributions can never cancel at an
individual order. The next section proves that they cannot cancel at every
large order. Higher Taylor coefficients of the prescribed curve enter lower
orders of the large-order expansion; no formula for them is needed here.
\end{remark}

\section{Tied poles and the local nonextension theorem}
\label{sec:obstruction}
\begin{lemma}[A finite oscillatory sum cannot disappear]
\label{lem:cesaro}
If $z_1,\dots,z_s$ are distinct complex numbers of modulus one and not all
$A_j$ vanish, then
\begin{equation}
 \lim_{N\to\infty}\frac1N\sum_{n=1}^N
 \left|\sum_{j=1}^s A_jz_j^n\right|^2
 =\sum_{j=1}^s|A_j|^2>0.
 \label{eq:mean-square}
\end{equation}
In particular the inner sum has modulus bounded below by a positive constant
along an infinite subsequence.
\end{lemma}
\begin{proof}
Expand the square. Every off-diagonal average is a finite geometric average
with ratio $z_j\overline{z_k}\ne1$, so tends to zero. The diagonal terms are
constant. The subsequence assertion follows immediately.
\end{proof}

\begin{proposition}[No analytic cancellation along a moving curve]
\label{prop:no-cancellation}
The function \eqref{eq:moving-function} has no holomorphic continuation
through $v=0$ in either of the following cases:
\begin{enumerate}[label=(\roman*)]
\item $d=2$ and $0<|q_0|<1$;
\item $d\ge3$ and $0<|q_0|<\e^{-d}$.
\end{enumerate}
The same conclusion holds after adding any function holomorphic near zero.
\end{proposition}
\begin{proof}
The set $J$ has at most two elements. For $d=2$, distinct $a_1,a_2\in J$
cannot satisfy $a_1^2=a_2^2$, because that would give $a_1=-a_2$, whereas
both have the same positive real part. For $d\ge3$, the assumed bound gives
$\Rea a>d$. Every nearest pole has $|\Ima a|\le\pi$, so
$|\arg a|<\pi/d$. Two distinct nearest poles therefore have distinct $d$th
powers: their argument difference is strictly between $-2\pi/d$ and $2\pi/d$.

In Theorem~\ref{thm:moving}, divide the expression in parentheses by
$\rho^{-dn}$. Its principal part is a finite sum
\[
 \sum_{a\in J} A_a z_a^n,\qquad z_a=(\rho/a)^d,
\]
where $A_a=\gamma_a\e^{a h_1/2}/a$ for $d=2$ and
$A_a=\gamma_a/a$ for $d\ge3$. Every $A_a$ is nonzero and the $z_a$ are
distinct. Lemma~\ref{lem:cesaro} and the $o(1)$ error imply, for infinitely
many $n$,
\begin{equation}
 |f_n|\ge c\frac{(dn)!}{n!\rho^{dn}}
 \label{eq:subsequence-growth}
\end{equation}
with a fixed $c>0$. Stirling's formula shows that the $n$th root of the
right-hand side is asymptotic, up to the factor $c^{1/n}$, to
\[
 d^d\e^{1-d}\rho^{-d}n^{d-1},
\]
which diverges. A holomorphic germ has geometrically bounded Taylor
coefficients. Its Taylor coefficients would have to equal the boundary
coefficients $f_n$, a contradiction. Addition of a holomorphic function
changes the coefficients by at most a geometric sequence and cannot remove
\eqref{eq:subsequence-growth}.
\end{proof}

\begin{corollary}[Rational-point obstruction for partial theta level sets]
\label{cor:rational}
Let $u_0=2\pi i a/b$, and let $q(u)$ be holomorphic near $u_0$, with
$q(u_0)\ne0$. Assume $|q(u_0)|<1$ for $d=2$, or $|q(u_0)|<\e^{-d}$ for
$d\ge3$. Then the left-hand function
\[
 u\longmapsto\Theta_d(q(u),u)
\]
cannot extend holomorphically through $u_0$.
\end{corollary}
\begin{proof}
Put $v=u-u_0$, $q_0=q(u_0)$, and choose the local logarithm
$h(v)=\log(q(u_0+v)/q_0)$. The periodic coefficients are
$c_j=\e^{j^d u_0}$. Apply Proposition~\ref{prop:no-cancellation}.
\end{proof}
The small bound $\e^{-d}$ in the higher-degree case is a convenient sufficient
condition that avoids resonances between nearest-pole $d$th powers. It is not
claimed to be an optimal bound. No such extra smallness is needed for $d=2$.

\section{Natural boundaries of nonlinear implicit solutions}
\label{sec:natural}
\begin{theorem}[Analytic level sets inherit a natural boundary]
\label{thm:forcing}
Let $f$ be holomorphic near zero with $f(0)=0$ and $f\not\equiv0$. For each
integer $d\ge2$, the small solution $Q_f$ of \eqref{eq:forced} has an entire
imaginary diameter as a natural boundary, after the radius is made
sufficiently small.
\end{theorem}
\begin{proof}
Use Lemma~\ref{lem:root}, shrinking the radius so that $|Q_f|<1$ for $d=2$
and $|Q_f|<\e^{-d}$ for $d\ge3$. At every rational imaginary boundary point
$u_0$ for which $f(u_0)\ne0$, the root $Q_f(u_0)$ is nonzero.
If $Q_f$ extended holomorphically through this point, Corollary~\ref{cor:rational}
would say that $\Theta_d(Q_f(u),u)$ cannot do so. But that function equals
the holomorphic function $f(u)$ on the left. This is a contradiction.

The exceptional zeros of $f$ are isolated. Rational imaginary points not among
those zeros are dense in the diameter. If continuation existed through any
point of the diameter, its open neighborhood would contain one of these
already forbidden rational points. Restricting the same continuation to a
smaller neighborhood of that rational point gives the final contradiction.
\end{proof}
This proves Theorem~\ref{thm:main}(i)--(ii). In particular, the boundary is not
merely a sequence of inaccessible rational points: density and openness of
holomorphic continuation give nonextension at every point of the arc.

\subsection{Finite-core robustness and periodic amplitudes}
\begin{theorem}[Robustness under analytic finite-core changes]
\label{thm:robust}
Let $(b_j)$ be a nonzero periodic sequence, $J_0\ge1$, and suppose
\begin{equation}
 K(q,u)=A(q,u)+\sum_{j\ge J_0}b_jq^j\e^{j^d u},
 \label{eq:robust-kernel}
\end{equation}
where $A$ is holomorphic near $(0,0)$, $K(0,u)=0$, and
$\partial_qK(0,0)\ne0$. If $f$ is as in Theorem~\ref{thm:forcing}, the small
solution of $K(q,u)=f(u)$ is holomorphic in a left half-disc, smooth to its
diameter, and has that diameter as a natural boundary.
\end{theorem}
\begin{proof}
The uniform geometric derivative bounds used in Lemma~\ref{lem:root} still
hold, with a constant accounting for $\max|b_j|$ and the analytic term. The
nonzero linear derivative supplies the small implicit root, either by a
normalized contraction or by its usual uniform proof.

At a rational imaginary point, the product
$b_j\e^{j^d u_0}$ is a nonzero periodic sequence. Were the root analytic
there, the omitted finitely many terms and the composition $A(q(u),u)$ would
be holomorphic. Adding back those finitely many terms would contradict
Proposition~\ref{prop:no-cancellation}. The same density argument completes
the proof.
\end{proof}

\begin{corollary}[Changing finitely many slopes does not remove the boundary]
\label{cor:finite-slopes}
If the slopes in \eqref{eq:general-feedback} equal $j^d$ for all sufficiently
large $j$, and all the remaining slopes are finite, the canonical small inverse
has the same natural-boundary conclusion. Nonnegativity is not needed for this
boundary assertion, although it is convenient for the summability statements.
\end{corollary}
\begin{proof}
The difference between the altered kernel and $\Theta_d$ is a finite sum of
entire functions of $(q,u)$. The altered kernel still has linear
$q$-derivative equal to one at $(0,0)$.
Apply Theorem~\ref{thm:robust}.
\end{proof}

\section{The forward boundary and its geometry}
\label{sec:geometry}
\begin{proof}[Proof of Theorem~\ref{thm:main}(iii)--(iv)]
The derivative $Q_d'(u)$ extends continuously to zero and equals one there.
On a sufficiently small closed half-disc, $|Q_d'-1|\le1/2$. For two points
$u,v$ in that convex set, integration along their segment gives
\[
 |Q_d(v)-Q_d(u)-(v-u)|\le\tfrac12|v-u|.
\]
Thus $Q_d$ is injective, with a continuous inverse on its image, and is
biholomorphic in the interior. The arc $\Gamma_d(t)=Q_d(it)$ has nonzero
velocity $iQ_d'(it)$.

Suppose $U_d=Q_d^{-1}$ extended holomorphically through
$q_0=\Gamma_d(t_0)$. Its interior derivative tends to $1/Q_d'(it_0)\ne0$,
so the extension has nonzero derivative at $q_0$. The holomorphic inverse
function theorem would extend $Q_d$ through $it_0$, contradicting
Theorem~\ref{thm:forcing}. Hence the whole arc is a natural boundary for the
forward realization.

The finite Taylor expansion \eqref{eq:first-Q} is valid along the diameter
by the proved smoothness. Substitution $u=it$ gives
\eqref{eq:boundary-shape-main}. For $d=2$,
\begin{equation}
 \Gamma_2(t)=it+2t^2-\tfrac{i}{2}t^3-\tfrac23t^4
                     -\tfrac{29i}{8}t^5+O(t^6).
 \label{eq:boundary-five}
\end{equation}
The inverse's singular diameter approaches directions $\pm\pi/2$ exactly;
the forward singular arc approaches them with a positive real displacement
$2t^2+O(t^4)$. Every sector about the negative axis of opening greater than
$\pi$ contains sufficiently small nonzero points of these boundary arcs.
An extension to such a sector would contradict their nonextension.
\end{proof}

\begin{remark}[A geometric rather than algebraic obstruction]
All finite-order formal inversions remain valid. The obstruction appears
when infinitely many exponential feedback terms are realized analytically.
A smooth boundary graph may therefore be visible to arbitrarily long finite
jets while remaining impossible to cross holomorphically.
\end{remark}

\section{Convergent exponential blocks and uniform remainders}
\label{sec:blocks}
Natural boundaries do not prevent accurate evaluation on the permitted side.
We give a constructive realization which also connects the theorem back to
transseries. Set $P_d(u)=\e^uQ_d(u)$. Then
\begin{equation}
 P_d=u-\sum_{j\ge2}P_d^j\e^{(j^d-j)u}.
 \label{eq:P}
\end{equation}
Introduce a temporary variable $t$ in place of the first $u$ on the right.
The solution has a block expansion
\begin{equation}
 P_d(u)=\sum_{k\ge1}u^k H_{d,k}(u),
 \qquad H_{d,k}(u)=\sum_\beta h_{d,k,\beta}\e^{\beta u},
 \label{eq:blocks}
\end{equation}
where each inner sum is finite and each coefficient is an integer.

\begin{proposition}[Tree support and a convergent block certificate]
\label{prop:blocks}
Let
\begin{equation}
 C(t)=\frac{1+t-\sqrt{1-6t+t^2}}4=\sum_{k\ge1}c_kt^k,
 \qquad \rho_* =3-2\sqrt2.
 \label{eq:schroeder}
\end{equation}
Then
\begin{equation}
 0\le\beta\le k^d-k,
 \qquad \sum_\beta|h_{d,k,\beta}|\le c_k\le\rho_*^{-k}.
 \label{eq:block-bounds}
\end{equation}
The block expansion converges normally on every closed left half-disc
$|u|\le r_0<\rho_*$. Its tail after $k=K$ has modulus at most
\begin{equation}
 \frac{(|u|/\rho_*)^{K+1}}{1-|u|/\rho_*}.
 \label{eq:block-tail}
\end{equation}
\end{proposition}
\begin{proof}
Expand \eqref{eq:P} recursively as plane rooted trees whose internal vertices
have arity $j\ge2$. A tree with $k$ leaves contributes a sign, and the
frequency is the sum of $j^d-j$ over internal vertices. If their arities are
$j_v$, then $\sum_v(j_v-1)=k-1$. The polynomial
$(s+1)^d-(s+1)$ has nonnegative coefficients and vanishes at $s=0$;
its value at a sum of nonnegative arguments is at least the sum of its values.
Hence
\[
 \sum_v(j_v^d-j_v)\le k^d-k.
\]
All frequencies are nonnegative. Forgetting signs counts these trees by
$C=t+C^2/(1-C)$, which yields \eqref{eq:schroeder}. Since its coefficients
are nonnegative and $C(\rho_*)<1$, one has $c_k\le\rho_*^{-k}$.

For $\Rea u\le0$, each exponential has modulus at most one. The geometric
majorant proves normal convergence and \eqref{eq:block-tail}. The summed
series satisfies $|P_d(u)|\le C(|u|)<1$. The tree majorant also bounds
all the substitutions in \eqref{eq:P} by $\sum_{j\ge2}C(|u|)^j<\infty$.
Absolute convergence therefore verifies the identity throughout the stated
half-disc. The resulting $Q_d=\e^{-u}P_d$ equals the small branch near zero
and provides its holomorphic continuation to that half-disc. In particular
$|Q_d|\le\e^{\rho_*}C(\rho_*)<1$, so the original kernel remains
absolutely convergent there as well.
\end{proof}

\begin{proposition}[All-orders boundary-uniform remainder]
\label{prop:remainder}
Write the formal Taylor series of $P_d$ as $\sum_{n\ge1}a_{d,n}u^n$.
For $\Rea u\le0$, $r=|u|\le r_0<\rho_*$, and $N\ge2$,
\begin{equation}
 \left|P_d(u)-\sum_{n<N}a_{d,n}u^n\right|
 \le r^N M_{d,N}(r_0),
 \label{eq:remainder-bound}
\end{equation}
where
\begin{equation}
 M_{d,N}(r_0)=\frac{\rho_*^{-N}}{1-r_0/\rho_*}
 +\sum_{k=1}^{N-1}c_k\frac{(k^d-k)^{N-k}}{(N-k)!}.
 \label{eq:majorant}
\end{equation}
For the approximation to $Q_d$ obtained by multiplying the truncated
polynomial by $\e^{-u}$, multiply this bound by $\e^r$.
\end{proposition}
\begin{proof}
For $\Rea z\le0$, the integral Taylor remainder gives
\[
 \left|\e^z-\sum_{m<L}\frac{z^m}{m!}\right|\le\frac{|z|^L}{L!}.
\]
In block $k<N$, take $L=N-k$ and use \eqref{eq:block-bounds}. Blocks $k\ge N$
are bounded without expansion. These two sums give \eqref{eq:majorant}.
At any fixed Taylor degree only finitely many blocks contribute, so the
truncated coefficients are exactly $a_{d,n}$.
\end{proof}
For $d=2$ this is the predecessor's quadratic certificate \cite{repo-fine},
now recovered from the tree formulation. Its validity on the imaginary
boundary is essential; estimates only on proper left sectors would not give
the same critical-direction conclusion.

\section{An explicit optimized scale for every integer degree}
\label{sec:optimal}
\begin{theorem}[Sharp growth of the positive remainder majorant]
\label{thm:majorant}
For fixed $d\ge2$ and $r_0<\rho_*$,
\begin{align}
 \log M_{d,N}(r_0)
 ={}&(d-1)N\log N-dN\log\log N\\
 &+\left[d\log\frac{d}{d-1}-(d-1)\right]N+o(N).
 \label{eq:majorant-asymp}
\end{align}
In particular, for every $H>0$ there is $C_H$ such that
\begin{equation}
 M_{d,N}(r_0)\le C_H H^N(N!)^{d-1}\qquad(N\ge2).
 \label{eq:zero-type}
\end{equation}
\end{theorem}
\begin{proof}
For $2\le k<N$, let $T_{N,k}=c_k(k^d-k)^{N-k}/(N-k)!$.
The logarithm of a sum of at most $N$ positive terms differs from its maximum
by at most $\log N$. Put $t=k/N$. The bounds $1\le c_k\le\rho_*^{-k}$,
$k^d-k\le k^d$, and Stirling's formula, with error uniform in $N-k\ge1$,
give
\[
 \frac1N\log T_{N,k}
 \le(1-t)\{(d-1)\log N+d\log t-\log(1-t)+1\}
       +t\log(\rho_*^{-1})+O(\log N/N).
\]
Set $a=t\log N$. After subtracting
$(d-1)\log N-d\log\log N$, the right-hand side becomes
\[
 d\log a+1-(d-1)a+R_d(t)+O(\log N/N),
\]
where
\[
 R_d(t)=-dt\log t-(1-t)\log(1-t)-t+t\log(\rho_*^{-1}).
\]
The function $R_d$ is bounded and tends to zero as $t\downarrow0$.
Since $d\log a+1-(d-1)a$ tends to minus infinity at both endpoints of
$(0,\infty)$, maximizing indices must have $a$ in a fixed compact subinterval.
There $t\to0$ uniformly. The limiting maximum is attained at
$a=d/(d-1)$ and equals
$d\log(d/(d-1))-(d-1)$.

For a matching lower bound, take
$k=\lfloor dN/((d-1)\log N)\rfloor$ and use $c_k\ge1$.
Replacing $k^d-k$ by $k^d$ changes the logarithm by
$(N-k)\log(1-k^{1-d})=o(N)$, including $d=2$.
The same calculation attains the limiting maximum. The geometric term in
\eqref{eq:majorant} has logarithm $O(N)$ and is smaller. This proves
\eqref{eq:majorant-asymp}. Subtracting
$(d-1)\log(N!)$ gives $-dN\log\log N+O(N)$, which implies
\eqref{eq:zero-type}.
\end{proof}

\begin{corollary}[Certified optimized truncation]
\label{cor:optimal}
Set $s=d-1$, $L=\log(1/r)$, and
\begin{equation}
 N_d(r)=\left\lfloor d^{-d/s}r^{-1/s}L^{d/s}\right\rfloor.
 \label{eq:optimal-index}
\end{equation}
For every $\varepsilon>0$, all sufficiently small $r>0$, and uniformly on
$|u|=r$, $\Rea u\le0$,
\begin{equation}
 \left|Q_d(u)-\e^{-u}\sum_{n<N_d(r)}a_{d,n}u^n\right|
 \le\exp\!\left[-\left(s d^{-d/s}-\varepsilon\right)
                     r^{-1/s}L^{d/s}\right].
 \label{eq:optimized}
\end{equation}
\end{corollary}
\begin{proof}
Insert \eqref{eq:optimal-index} in \eqref{eq:majorant-asymp}. One has
\[
 \log N=\frac{L+d\log L-d\log d}{s}+o(1),\qquad
 \log\log N=\log L-\log s+o(1).
\]
It follows that $\log(r^N M_{d,N})=-sN(1+o(1))$.
The factor $\e^r$ in Proposition~\ref{prop:remainder} is negligible at this
scale. The arbitrary $\varepsilon$ absorbs the remaining lower-order terms.
\end{proof}
The constant is $1/4$ for $d=2$ and $2/(3\sqrt3)$ for $d=3$.
It is a sharp constant for the leading optimization of this positive
majorant. A matching lower bound for the \emph{actual signed truncation error}
is not proved. In the transseries coordinate $u=-\e^{-X}$,
\eqref{eq:optimized} is a bound on the scale
\[
 \exp\!\left[-\left((d-1)d^{-d/(d-1)}-o(1)\right)
       \e^{X/(d-1)}X^{d/(d-1)}\right].
\]
This upper bound alone does not establish a nonzero additional transseries
sector on that scale.

\section{Fine summability and the angular obstruction}
\label{sec:summability}
\subsection{Conventions and the classical implication used}
For $\widehat f(z)=\sum_{n\ge0}f_nz^n$, use the integrated Borel transform
\begin{equation}
 \B\widehat f(\zeta)=\sum_{n\ge0}\frac{f_n}{n!}\zeta^n.
 \label{eq:Borel}
\end{equation}
A series is \emph{finely summable in direction $\pi$} if this germ continues
to a tube of some fixed width about $(-\infty,0]$, with a bound
$C\e^{H|\zeta|}$ there. The sum on $z=-r$ is
\begin{equation}
 \mathcal S_\pi\widehat f(-r)=\frac1r\int_0^\infty
          \e^{-s/r}\B\widehat f(-s)\dd s.
 \label{eq:sum}
\end{equation}
This is the fine convention of Sauzin's Sections 7 and 9, reflected to the
negative ray \cite{sauzin}. By \emph{angular $1$-summability near $\pi$} we
mean continuation to an open Borel sector containing that direction, with
exponential bounds uniform on each smaller angular sector. Contour rotation
then supplies a sum in a physical sector of opening greater than $\pi$.

The Nevanlinna--Sokal theorem says that holomorphy on a tangent disc
$\Rea(-1/z)>R^{-1}$, together with uniform remainders bounded by
$CA^NN!|z|^N$, reconstructs the function by fine Borel summation
\cite{sokal}. We use this classical theorem, not a replacement theorem.

\subsection{Recovery of the positive quadratic result}
For $d=2$, \eqref{eq:zero-type} and Proposition~\ref{prop:remainder} imply
uniform Gevrey--1 remainders for $P_2$ on a closed left half-disc. Every
sufficiently small tangent disc lies there. Nevanlinna--Sokal therefore
identifies $P_2$ as its fine sum. Multiplication by $\e^{-u}$ preserves that
class and identifies $Q_2$ as the fine sum of $\Qhat_2$.
Taking a boundary coefficient limit in the all-orders estimate shows that
its coefficients satisfy $|[u^n]P_2|\le C_HH^nn!$ for every $H>0$.
Hence $\B\widehat P_2$ is entire. The same holds after multiplication by
$\e^{-u}$.

For completeness, the nonlinear closure needed to pass to $U_2$ can be
seen by conjugating a tangent-to-identity series at zero to one at infinity.
If $f(z)=z+c+\varphi(z)$, $\varphi=O(z^{-1})$, has a fine-summable Borel
minor $\phi$, first remove the constant translation. The inverse correction
has the formal representation
\begin{equation}
 \chi=\sum_{k\ge1}\frac{(-1)^k}{k!}
                   \partial_z^{k-1}(\varphi^k),\qquad
 \mathscr B\chi(\xi)=-\sum_{k\ge1}\frac{\xi^{k-1}}{k!}\phi^{*k}(\xi).
 \label{eq:inverse-minor}
\end{equation}
If $|\phi(\xi)|\le M\e^{c_0|\xi|}$ on a convex tube, convolution along a
straight simplex gives
\[
 |\phi^{*k}(\xi)|\le M^k\e^{c_0|\xi|}
                      \frac{|\xi|^{k-1}}{(k-1)!}.
\]
Thus \eqref{eq:inverse-minor} converges normally and is bounded by
$M\exp((c_0+2\sqrt M)|\xi|)$. Termwise Laplace integration and the
convergent analytic Lagrange formula at sufficiently large $\Rea z$
identify its sum as the actual inverse. Translations, reciprocals with
nonzero constant term, and multiplication or division by $z$ preserve the
fine class, after narrowing the tube for differentiation. Conjugation by
$z\mapsto-1/z$ therefore proves the claim for $U_2$. These are the
half-strip operations also used in the predecessor \cite{repo-fine}; the
angular version is classical \cite{sauzin}.

If the input minors are entire, the same estimates on each fixed compact
disc, with constants depending on that disc, prove that the transformed
minors are entire. Consequently $\B\Uhat_2$ is entire as well. These local
compact-disc bounds make no assertion about exponential growth on an
unbounded angular sector.

\subsection{The negative answer to angular summability}
\begin{proof}[Proof of Theorem~\ref{thm:quadratic-main} for the exact quadratic model]
The preceding subsection supplies fine summability, entire Borel transforms,
and identification of the actual sums. Suppose $\Qhat_2$ were angularly
$1$-summable near $\pi$. Its angular sum agrees with its fine sum along the
negative ray, because both use the same continued Borel germ in
\eqref{eq:sum}. Contour rotation makes the angular sum holomorphic in a
sector about the negative axis of opening greater than $\pi$. The identity
theorem identifies it with $Q_2$ on their left-hand overlap. It would thus
continue $Q_2$ through small points of its imaginary natural boundary,
contradicting Theorem~\ref{thm:main}.

If $\Uhat_2$ were angularly summable, the classical closure of angularly
summable tangent-to-identity series under inversion would make $\Qhat_2$
angularly summable, which has just been excluded. Alternatively, its fine
sum is already $U_2$, and the same sector argument meets the natural boundary
$\Gamma_2$ directly.
\end{proof}

\begin{corollary}[Entire Borel transforms with no angular exponential bound]
\label{cor:Borel-growth}
For either $B=\B\Qhat_2$ or $B=\B\Uhat_2$, for every $\varepsilon>0$ and
every $C,H>0$, there are arbitrarily large $\zeta$ with
$|\arg\zeta-\pi|<\varepsilon$ (using the argument centered at $\pi$) such that
\[
 |B(\zeta)|>C\e^{H|\zeta|}.
\]
Thus the obstruction is growth at Borel infinity, not a finite singularity
of these entire Borel transforms.
\end{corollary}
\begin{proof}
A uniform exponential bound on such a sector, together with entireness,
would give angular summability. If the bound held beyond some radius, boundedness on the remaining compact
disc would give a global bound after increasing $C$. That too would give
angular summability. Thus violations occur at arbitrarily large points.
\end{proof}
The conclusion does not classify growth on each individual nonreal ray.
Failure of a sector-wide bound is stronger than a failure at one fixed
finite point, but weaker than a pointwise growth law on every ray.

\subsection{Completion for finitely modified quadratic slopes}
Normalize the small inverse by $P=\e^{\lambda_1u}Q$. Its tree frequency at an
internal vertex of arity $j$ is $\lambda_j-j\lambda_1$.
For slopes equal to $j^2$ eventually, there is a constant $C$ such that a
$k$-leaf tree has frequency $\beta$ satisfying
\[
 -Ck\le\beta\le Ck^2.
\]
Indeed $\sum(j-1)=k-1$, the number of internal vertices is at most $k-1$,
and only finitely many arities differ from the quadratic rule.
On $|u|\le r_0$, $\Rea u\le0$, negative frequencies cost at most
$\e^{Ckr_0}$. This merely changes the geometric tree-count constant.
The proof of Proposition~\ref{prop:remainder} now uses $(Ck^2)^{N-k}$ and
an adjusted geometric constant. The saddle estimate still gives
\[
 \log M_N=N\log N-2N\log\log N+O(N),
\]
so the remainders are Gevrey--1 of zero type. The same summation and inverse
arguments apply. Corollary~\ref{cor:finite-slopes} supplies the natural
boundary, hence excludes angular summability. This proves the remaining
assertion of Theorem~\ref{thm:quadratic-main}.

\section{Further extensions and sharp limitations}
\label{sec:extensions}
\subsection{Affine-quadratic exponents}
The natural-boundary mechanism also treats
\[
 \lambda_j=A j^2+Bj+C,\qquad A>0,
\]
with real $B,C$, near the origin. Set $w=Au$ and $z=q\e^{Bu}$.
Then
\[
 \sum_{j\ge1}q^j\e^{\lambda_j u}
 =\e^{Cu}\Theta_2(z,w).
\]
The level equation becomes
$\Theta_2(z,w)=(w/A)\e^{-Cw/A}$, an analytic nonzero forcing to which
Theorem~\ref{thm:forcing} applies. Multiplication of the moving root by a
nonvanishing holomorphic factor and the positive linear change of $u$ cannot
remove the boundary. Finite perturbations are again permitted.

\subsection{The essential small-parameter hypothesis}
The local theorem assumes $q_0\ne0$. At $q_0=0$ the entire kernel vanishes,
and the holomorphic curve $q\equiv0$ is an obvious exception. The natural
boundary proof does not overlook it: it proves nonextension on a dense set
of nonzero-root points and then uses openness of continuation to reach the
isolated zero points, including the origin.

For $d\ge3$, the smallness hypothesis separates the $d$th powers of nearest
poles. Without it, cancellation of a leading nearest-pole pair is a real
possibility in the asymptotic formula. One must then inspect further orders
or further poles. The present theorem does not silently assume that such a
pair cannot occur.

\subsection{No contradiction with local resurgence of other theta expansions}
The partial-theta literature studies several different expansions, including
ones at rational boundary points with periodic amplitudes, and their Borel
singularities \cite{hlss}. Our coefficient asymptotics concern a geometric
weight $q_0^j$ and a holomorphically moving $q$, and the nonlinear conclusion
concerns the inverse near the vanishing level. Neither the expansion point
nor the summation problem may be exchanged without proof.

In particular, an entire Borel transform is endlessly continuable in the
finite plane in the simplest sense, while still failing the growth
requirements for angular Laplace summation. The word ``resurgent'' alone
therefore does not answer the angular question settled here.

\section{Verification, reproducibility, and proof audit}
\label{sec:verification}
\subsection{Exact finite checks}
The supplied \path{code/verify.py} uses rational arithmetic for three
independent constructions: the original feedback recurrence, formal
compositional inversion, and the exponential-block tree recurrence.
At the recorded default order $24$, for $d=2,3,4$, it performed
\textbf{324 exact comparisons}. Both inverse compositions, the inverse
kernel identity, the displayed low coefficients, block support, block mass,
and block-to-Taylor conversion passed.

The block calculations use blocks through degree $10$. The ordinary formal
series calculations use degree $24$. These are deliberately separate
checks, not a single recurrence compared with itself.
The CSV files retain exact rational coefficients. None of these finite
checks proves a natural boundary or an infinite-order asymptotic theorem.

\subsection{High-precision diagnostics}
The optional diagnostic section ran at $110$ decimal digits with mpmath.
It checked the moving-curve asymptotics at ten $(d,n)$ pairs, four rational
boundary roots, $352$ quadratic Fourier magnitudes, and a tied-pole
mean-square example. These are floating-point diagnostics, not interval
certificates.

For $q(v)=q_0\exp(3v/7+v^2/9)$ with constant periodic coefficient $c_j=1$,
the ratio of the exact finite coefficient formula to its predicted leading
term was as follows. The quadratic case uses $q_0=1/5$; the cubic case uses
$q_0=1/30$.
\begin{center}
\begin{tabular}{rcc}
\toprule
$n$ & $d=2$ ratio & $d=3$ ratio\\
\midrule
8  & 1.003554662 & 1.075228530\\
16 & 1.001774139 & 1.035862320\\
32 & 1.000884142 & 1.017556494\\
64 & 1.000441140 & 1.008690883\\
80 & 1.000352752 & 1.006939094\\
\bottomrule
\end{tabular}
\end{center}
These values are consistent with the limits; the proof of the limits is
Section~\ref{sec:moving}, not this table.

At $u=2\pi i/b$, the boundary equation is finite and rational in $q$:
\begin{equation}
 \Theta_2(q,2\pi i/b)
 =\frac{\sum_{j=1}^{b}\e^{2\pi i j^2/b}q^j}{1-q^b}.
 \label{eq:boundary-rational}
\end{equation}
Solving it for the small root gives the following diagnostic for the
quadratic displacement in \eqref{eq:boundary-five}:
\begin{center}
\begin{tabular}{rcc}
\toprule
$b$ & $t=2\pi/b$ & $\Rea Q_2(it)/t^2$\\
\midrule
127  & 0.04947390006 & 1.998438396\\
257  & 0.02444819186 & 1.999605886\\
509  & 0.01234417546 & 1.999898701\\
1021 & 0.00615395231 & 1.999974770\\
\bottomrule
\end{tabular}
\end{center}
The maximum high-precision residual was below $1.4\cdot10^{-114}$.
This is a residual in the stated floating-point computation, not a certified
error bound of that magnitude for every printed digit.

\subsection{The proof's load-bearing points}
The natural-boundary theorem rests on four analytic facts proved above:
the pole expansion of the moment generating function; a global factorial
majorant valid even near the end of the coefficient sum; a noncancellation
argument when two nearest poles tie; and a contradiction obtained only
\emph{after assuming} that the implicit root is holomorphic across a point.
A fixed-parameter gap theorem, numerical derivative growth, or a plot of
boundary values would not replace any of these steps.

The positive quadratic summation result additionally uses the classical
Nevanlinna--Sokal theorem and the standard convolution/Lagrange inverse
mechanism. The new natural-boundary proof itself does not depend on a
summability theorem, the accuracy of the numerical tables, or an unverified
claim about the larger transseries field.

\subsection{Suggested formalization layers}
A realistic formalization would separate finite Fourier inversion and finite
coefficient identities from complex analytic estimates. First formalize
\eqref{eq:exact-moving}, \eqref{eq:factorial-ratio}, and
\eqref{eq:double-majorant}. Next package the Cauchy principal-part estimate
and dominated-convergence argument. Only then assemble the periodic boundary
obstruction and its implicit-function corollary. The summability corollary is
a later, independent interface. No Lean source is included or represented as
checked in this package.

\section{Further research questions and topics}
\label{sec:questions}
The following problems are not used as premises in the proofs.

\subsection{Boundary jets of the actual implicit solution}
\begin{question}
At a fixed nonzero rational imaginary point $u_0$, determine the precise
large-order asymptotics of the actual boundary derivatives
$Q_d^{(n)}(u_0)/n!$, including constants and cancellations.
\end{question}
The moving-curve theorem assumes a holomorphic curve and proves that such a
curve cannot be the implicit root. It does not give the root's own
large-order law. A nonlinear asymptotic inversion of the periodic moment
system is the next step. Establishing exact boundary Gevrey order would
quantify how far the $C^\infty$ boundary is from analyticity.

\subsection{The angular growth of the entire Borel transform}
\begin{question}
For $B=\B\Qhat_2$ or $\B\Uhat_2$, determine the growth as
$|\zeta|\to\infty$ on rays adjacent to the negative axis, and identify the
largest curvilinear regions carrying exponential bounds.
\end{question}
Corollary~\ref{cor:Borel-growth} rules out every fixed angular neighborhood,
but leaves room for widening tubes and parabolic regions. An explicit
transition law could explain how fine summation approaches its angular
failure and improve numerical contour choices.
% ed. (2026-09-29): a partial answer in an earlier-filed package (batch 49).
\begin{ednote}
For $B=\B\Qhat_2$ the maximum modulus on circles is known:
\path{Signed_Quadratic_Feedback_Inversion/} \cite{ed:sqi},
Theorem~\texttt{thm:borel}, gives $\max_{|\zeta|=R}|B(\zeta)|\sim-B(R)$, with
an explicit equivalent, which bounds the growth on every ray. The growth on
individual rays adjacent to the negative axis is not determined there.
\end{ednote}

\subsection{Higher-degree pole collisions}
\begin{question}
Remove the convenient bound $|q_0|<\e^{-d}$ from
Proposition~\ref{prop:no-cancellation} for $d\ge3$, or classify the exceptional
curves on which leading $d$th-power pole collisions cancel.
\end{question}
The exact obstruction is equality of nearest-pole $d$th powers, not a
failure of the coefficient identity. Subleading terms and additional pole
circles should be analyzed before asserting a general nonextension theorem
for the whole $q$-disc.

\subsection{General polynomial and nonpolynomial slopes}
\begin{question}
Extend the natural-boundary theorem from $j^d$ and affine quadratics to
arbitrary positive-leading-coefficient integer polynomials, and then to
nonintegral powers or asymptotically polynomial slopes.
\end{question}
For integer polynomials, rational boundary phases remain periodic, but the
lower polynomial terms alter the moment asymptotics. For noninteger powers,
periodicity is lost; a different replacement for the finite Fourier pole
spectrum is needed. Merely knowing coefficient Gevrey growth does not settle
analytic continuation.

\subsection{Actual optimal errors rather than positive majorants}
\begin{question}
Does the signed error for optimally truncated $Q_d$ attain the scale and
constant in \eqref{eq:optimized} along some nonexceptional ray? If not, what
is its true leading scale?
\end{question}
A positive tree majorant discards cancellations. Matching lower bounds
would require a direct analysis of the signed blocks, not another
optimization of the same upper estimate.
% ed. (2026-09-29): the least formal term at d = 2 is determined by a package filed beside this one (batch 50).
\begin{ednote}
For $d=2$ the least formal term of $\Qhat_2$ on the positive axis is
determined by \path{Signed_Condensation_Sharp_Quadratic_Reversion/}
\cite{ed:scs}, Theorem~\texttt{thm:least}: at $a=1$ it is
$\exp(-(1/4+o(1))\log^2(1/t)/t)$, with the constant $1/4$ of
\eqref{eq:optimized} at $d=2$. That is a statement about formal terms; it
gives no lower bound for the analytic truncation error, so the question
stays open.
\end{ednote}

\subsection{Global boundary geometry and inverse branches}
\begin{question}
Continue the canonical root along the left half-plane as far as possible.
Locate any interior critical points or poles and determine the global image
of the imaginary boundary under $Q_d$.
\end{question}
The current theorem is local near zero. A global branch may encounter
critical points before the geometric-series condition $|Q_d|<1$ fails.
The local natural boundary is proved; its global continuation and topology
are separate questions.

\subsection{Summation procedures for higher-degree feedback}
\begin{question}
For $d\ge3$, identify a summation procedure that reconstructs the canonical
literal branch from its Gevrey--$(d-1)$, logarithmically improved series,
without assuming an angular continuation which the branch cannot have.
\end{question}
The closed-half-disc remainder and convergent block expansion already give
an evaluation procedure. A generalized fine summation or acceleration
theorem would provide a transform-side characterization. It must distinguish
an arbitrary function with the same formal expansion from the branch
selected by the literal infinite feedback equation.

\subsection{Effective boundary certificates}
\begin{question}
Turn the moment and mean-square proof into a finite, certified procedure
that rules out a proposed analytic continuation radius or derivative bound.
\end{question}
At rational boundary points the pole set is explicit. Quantitative errors
in Theorem~\ref{thm:moving}, together with an effective lower bound for a
short mean-square average, would convert the qualitative contradiction into
machine-checkable analytic certificates.

\section{Conclusion}
The infinite feedback equation supplies a concrete separation between formal
inversion, smooth boundary realization, fine Borel reconstruction, and
angular analytic continuation. In the quadratic model, fine reconstruction
works and the Borel transforms are entire, yet every angular enlargement of
the negative summation direction fails. The obstruction survives finite-core
changes and analytic forcing because a moving holomorphic parameter cannot
remove the factorial moment contribution of the periodic boundary spectrum.

The same proof gives natural boundaries for all integer-power feedback
kernels $j^d$, while the tree expansion gives usable, uniform remainder
bounds on the permitted side. The result is therefore not simply an
obstruction: it identifies both what continuation cannot do and what
certified asymptotic evaluation can still accomplish.

\appendix
\section{Reproduction commands and package contents}
The article source is standalone and uses an embedded bibliography. From the
package directory, run
\begin{verbatim}
sh build.sh
python code/verify.py --order 24 --exact-only
python -m pip install -r requirements.txt
python code/verify.py --order 24
\end{verbatim}
The first verification command needs only the Python standard library. The
second full run additionally uses mpmath. Running the program replaces its
own data files; an exact-only run marks diagnostics as not run rather than
pretending they passed. No program accesses the network or changes the
upstream repository.

The package includes source and PDF, the verification program, exact
coefficient and block files, numerical diagnostic tables, a recorded run,
source-provenance notes, and a proof-status audit. PDF compilation and
rendering checks are recorded separately from mathematical verification.

\begin{thebibliography}{99}
\raggedright
\bibitem{repo}
V. Reshetnikov, \emph{ProveIt}, repository and transseries documentation.
Snapshot \texttt{ebd8344bca77d8352cd745b7df374618a290a029},
accessed 29 September 2026.
\url{https://github.com/VladimirReshetnikov/ProveIt}.

\bibitem{repo-reg}
ProveIt research document, \emph{A sharp regularity classification for
countable exponential-feedback transseries}, 2026.
Repository directory
\path{Analysis/Transseries/docs/series-and-transseries/Exponential_Feedback_Regularity_Classification/}.
Related background; the fine-summation predecessor supplies the comparison
used in this article.

\bibitem{repo-sector}
ProveIt research document, \emph{Sharp negative-direction summability for
countable exponential feedback: A coefficient-to-remainder principle,
optimal logarithmic scales, and sectorial inversion}, 29 September 2026.
Repository directory
\path{Analysis/Transseries/docs/series-and-transseries/Sharp_Negative_Direction_Summability_Exponential_Feedback/}.
Source blob \texttt{f82840478f1e9034b3de05f7c8d0ca5a612274b4}.

\bibitem{repo-fine}
ProveIt research document, \emph{Negative-Ray Summation of Countable
Exponential-Feedback Transseries: A sharp slope threshold, an inverse Bessel
representation, and logarithmic-Gevrey error certificates},
29 September 2026. Repository directory
\path{Analysis/Transseries/docs/series-and-transseries/Negative_Ray_Summation_Exponential_Feedback/}.
Source blob \texttt{eafc690678f7ca2cd75a8310f1a0889b19c7cdb5}.

% ed. (2026-09-29): bibliography entries for the editorial notes (packages filed beside this one).
\bibitem{ed:nbq}
\emph{Natural Boundaries of Quadratic Exponential Feedback}, ProveIt research
package filed in batch 49 (2026-09-29),
\path{Analysis/Transseries/docs/series-and-transseries/Natural_Boundaries_Quadratic_Exponential_Feedback/}.
Editorial addition.

\bibitem{ed:sqi}
\emph{Signed Quadratic-Feedback Inversion}, ProveIt research package filed in
batch 49 (2026-09-29),
\path{Analysis/Transseries/docs/series-and-transseries/Signed_Quadratic_Feedback_Inversion/}.
Editorial addition.

\bibitem{ed:scs}
\emph{Signed Condensation in Quadratic-Feedback Transseries}, ProveIt research
package filed in batch 50 (2026-09-29),
\path{Analysis/Transseries/docs/series-and-transseries/Signed_Condensation_Sharp_Quadratic_Reversion/}.
Editorial addition.

\bibitem{sokal}
A. D. Sokal, An improvement of Watson's theorem on Borel summability,
\emph{Journal of Mathematical Physics} \textbf{21} (1980), 261--263.
\href{https://doi.org/10.1063/1.524408}{doi:10.1063/1.524408}.

\bibitem{sauzin}
D. Sauzin, \emph{Introduction to 1-summability and resurgence},
arXiv:1405.0356, 2014. In particular Sections 7, 9, 13, and 17.
\url{https://arxiv.org/abs/1405.0356}.

\bibitem{hlss}
L. Han, Y. Li, D. Sauzin, and S. Sun,
\emph{Resurgence and Partial Theta Series},
arXiv:2112.15223, version 3, 7 July 2022.
\url{https://arxiv.org/abs/2112.15223}.

\bibitem{folsom}
A. Folsom, Asymptotic expansions, partial theta functions, and radial limit
differences of mock modular and modular forms,
\emph{International Journal of Number Theory} (2020).
\href{https://doi.org/10.1142/S1793042120400126}{doi:10.1142/S1793042120400126}.

\bibitem{nikolaev}
N. Nikolaev, Gevrey asymptotic implicit function theorem,
\emph{L'Enseignement Math\'ematique} \textbf{70} (2024), 251--282.
\href{https://doi.org/10.4171/LEM/1061}{doi:10.4171/LEM/1061}.
\end{thebibliography}
\end{document}
